with @@M0@@ a @@M1@@-sheaf and @@M2@@. Then
@@M3@@
but
@@M4@@
by induction on @@M5@@ and @@M6@@; hence (i) by induction on @@M7@@. Finally, by 1.7 and induction on @@M8@@, @@M9@@ and @@M10@@ is a @@M11@@-sheaf; hence (ii).

\sgaStableTarget{sga6.e13.def.1-9}\textbf{Definition 1.9.} \emph{The following polynomials with rational coefficients @@M13@@, defined by induction, are called the \emph{@@M12@@-polynomials}}~:
@@M14@@

\sgaStableTarget{sga6.e13.rem.1-10}\textbf{Remark 1.10.} For the @@M15@@-polynomials, one sees at once by induction on @@M16@@ that
@@M17@@

\sgaStableTarget{sga6.e13.thm.1-11}\textbf{Theorem 1.11.} \emph{(The main theorem on @@M18@@-sheaves). Let @@M19@@ be a @@M20@@-sheaf on @@M21@@ and @@M22@@ the Hilbert polynomial of @@M23@@; let @@M24@@ be a sequence of integers @@M25@@ and @@M26@@ the integer given by the @@M27@@-th @@M28@@-polynomial (1.9). Then @@M29@@ and @@M30@@ is @@M31@@-regular. In particular, if @@M32@@, then @@M33@@ is @@M34@@-regular.}

